\section{Introduction}\label{sec:introduction}

A Hecke eigensheaf realizes a local system on a curve through the geometry of
modifications of bundles. For a punctured curve, the local monodromy of the
parameter must also be reflected in the level structure at the puncture.
We study Borel level and tame regular-unipotent monodromy. The required
automorphic object is a locally constructible perverse sheaf, and its
eigenisomorphisms must hold over the entire space of Hecke legs, including
their collision diagonals.

\subsection{The setting and the main theorem}
Let $k=\overline{\mathbb F}_q$ have characteristic $p$, and let
$E=\overline{\mathbb Q}_\ell$, where $\ell\ne p$.
Let $X_0/\mathbb F_q$ be a smooth projective geometrically connected curve of
genus $g\ge2$, with a rational point $x_0$.
Write $X=X_0\times_{\mathbb F_q}k$, $x=(x_0)_k$, and $U=X\setminus\{x\}$.
Let $G_0/\mathbb F_q$ be split, simple, and simply connected of positive rank,
and fix a Borel subgroup $B_0\subset G_0$.
Put $G=(G_0)_k$, $B=(B_0)_k$, and $\mathfrak g=\Lie(G)$.
The dual group $\Gd/E$ is adjoint.

We make the following assumptions on the characteristic. In the statements
below, $M$ ranges over Levi subgroups of $G$, $\mathfrak m=\Lie(M)$,
$T_M\subset M$ is a maximal torus, and $W_M=N_M(T_M)/T_M$.
\begin{enumerate}[label=\textup{(H\arabic*)},ref=H\arabic*]
\item\label{hyp:form}
There is a nondegenerate $G$-invariant symmetric form $\kappa$ on
$\mathfrak g$, and its restriction to the Lie-algebra center
$\mathfrak z(\mathfrak m)$ is nondegenerate for every $M$.
\item\label{hyp:chevalley}
Restriction induces an isomorphism
\[
 k[\mathfrak m]^M\xrightarrow{\sim}k[\Lie(T_M)]^{W_M}
 \qquad\text{for every }M.
\]
\item\label{hyp:centralizer}
The scheme-theoretic centralizer in $G$ of every semisimple element of
$\mathfrak g$ is a Levi subgroup.
\item\label{hyp:nilpotent}
For every field extension $k'/k$, each nilpotent element of
$\mathfrak g\otimes_k k'$ belongs to $\Lie(R_u(P))$ for some parabolic
subgroup $P\subset G_{k'}$ defined over $k'$.
\end{enumerate}
Here semisimple means geometrically conjugate into a toral Lie algebra, and
nilpotent means that the geometric adjoint-orbit closure contains zero.
We impose no additional condition that $p$ be prime to the order of a Weyl
group.

A parameter $\sigma$ is a continuous homomorphism
\begin{equation}\label{eq:parameter}
 \rho:\pi_1^{\et}(U,\bar u)\longrightarrow\Gd(L),
\end{equation}
up to $\Gd(E)$-conjugacy, where $L/\mathbb Q_\ell$ is finite inside $E$.
We require its image to be Zariski dense. It defines a tensor local system
$V\mapsto V_\sigma$ for $V\in\Rep_E(\Gd)$.
We further require that $\rho$ kill wild inertia at $x$ and that its tame
inertia action factor through $\mathbb Z_\ell(1)$.
After choosing a compatible system of $\ell$-power roots of unity, let
$t_\ell:I_x\to\mathbb Z_\ell$ be the resulting quotient map.
The regular-unipotent condition is
\begin{equation}\label{eq:regular-monodromy}
 \rho(\gamma)=\exp\bigl(t_\ell(\gamma)N\bigr),\qquad \gamma\in I_x,
 \qquad \dim Z_{\Gd}(N)=\rk(\Gd),
\end{equation}
for a nilpotent $N\in\Lie(\Gd)(L)$ after a finite extension of $L$ and a
conjugation. The exponential is the finite nilpotent exponential.
Changing the chosen generator rescales $N$ by a unit, so the condition does
not add a generator or a logarithm to the automorphic data.
No Frobenius descent is assumed for $\sigma$.

Let
\[
 \cA=\Bun_{G,B,x}(X)
\]
be the stack of $G$-bundles $P$ on $X$ with a $B$-reduction $\beta$ at $x$.
A geometric $E$-adic sheaf $M$ on $\cA$ is \emph{locally constructible
perverse} if, for every smooth map $f:S\to\cA$ from a smooth finite-type
$k$-scheme of constant relative dimension $d$, the complex $f^*M[d]$ is
bounded constructible and perverse. This condition is local on the bundle
stack; it requires no bound on the Harder--Narasimhan strata supporting $M$.

Using $\kappa$, a cotangent vector to $\cA$ is represented by
\[
 (P,\beta,\phi),\qquad
 \phi\in H^0\bigl(X,\ad(P)\otimes\omega_X(x)\bigr),\qquad
 \Res_x(\phi)\in\Lie(R_u(B_\beta)).
\]
The \emph{parabolic global nilpotent cone} $\Lambda_{\mathrm{par}}$ consists
of these triples for which $\phi$ is generically nilpotent.
The singular-support condition below is interpreted by pulling this cone
back to the cotangent bundle of each smooth scheme chart.

For a finite set $I$ and $V_i\in\Rep_E(\Gd)$, let
$\Hecke_{I,(V_i)}$ denote the Hecke functor away from $x$, with target
sheaves on $\cA\times U^I$.
We use the relative Satake normalization: on the open Schubert cell of
dimension $d_\lambda$ the simple kernel is
$E[d_\lambda](d_\lambda/2)$, with no additional moving-leg shift.
These dimensions are even because $G$ is simply connected.
In particular, the tensor-unit functor sends $M$ to
$M\boxtimes E_{U^I}$.
Section~\ref{sec:foundations} fixes the corresponding tensor and fusion
conventions.

\begin{theorem}\label{thm:main}
Assume \ref{hyp:form}--\ref{hyp:nilpotent}. For every Zariski-dense
parameter \eqref{eq:parameter} satisfying \eqref{eq:regular-monodromy},
there is a nonzero locally constructible perverse geometric $E$-adic sheaf
$M$ on $\cA$ such that
\[
 \SS(M)\subset\Lambda_{\mathrm{par}}.
\]
For every finite set $I$ and every family $(V_i)_{i\in I}$, it carries
isomorphisms
\begin{equation}\label{eq:main-eigen}
 \Hecke_{I,(V_i)}(M)
 \simeq M\boxtimes\Bigl(\mathop{\boxtimes}_{i\in I}(V_i)_\sigma\Bigr)
 \quad\text{on }\cA\times U^I,
\end{equation}
natural in the representations and coherently compatible with the tensor
unit, convolution, permutations of legs, and fusion along every collision
diagonal.
\end{theorem}

\subsection{Historical context}
The eigensheaf formulation of geometric Langlands grew from Drinfeld's
rank-two construction and Laumon's geometric formulation and proposed
construction for general linear groups \cite{Drinfeld,Laumon}. Frenkel,
Gaitsgory, and Vilonen reduced the unramified $\mathrm{GL}_n$ existence
problem to a vanishing theorem \cite{FGV}; Gaitsgory proved that theorem
\cite{GaitsgoryVanishing}. For general reductive groups in characteristic
zero, Beilinson and Drinfeld constructed eigensheaves from opers by
quantizing the Hitchin system \cite[\S0.2]{BD}. The tensor structure and
fusion of geometric Hecke operators are furnished by geometric Satake
\cite[\S\S4--5,14]{MV}.

Ramified eigensheaves have also been studied through automorphic data with
prescribed level and character. Yun's rigidity framework constructs
eigenvalues from suitable rigid automorphic data
\cite[\S4]{YunRigidity}; the present problem starts instead with a prescribed
Zariski-dense local system. In the characteristic-zero de Rham setting,
F{\ae}rgeman constructs coherent eigensheaves for irreducible regular-singular
parameters using the factorizable spectral input specified in his paper
\cite[Theorem~1.4.1.1,\S1.3.7]{Faergeman}. In the version cited here, regular
holonomicity and generic perversity in the ramified setting are left as expectations
\cite[Remark~1.4.1.2]{Faergeman}. Our theorem concerns geometric adic
coefficients in positive characteristic, Borel level, and local
constructibility together with perversity.

The geometry of the global nilpotent cone was developed by Laumon,
Faltings, and Ginzburg \cite{LaumonNilpotent,Faltings,GinzburgNilpotent}.
Beilinson established the singular-support theory for constructible
etale sheaves, and Barrett extended it to the adic coefficient categories
used here \cite{Beilinson,Barrett}. Arinkin, Gaitsgory, Kazhdan, Raskin,
Rozenblyum, and Varshavsky developed the restricted spectral stack and
its action on sheaves with nilpotent singular support \cite{AGKRRV}.
Their Hecke detection argument uses a central cocharacter of the Levi
centralizing a semisimple Higgs value and an isolated cotangent
intersection \cite[\S\S20.5,20.7--20.8, Appendix~H]{AGKRRV}.
The cotangent calculations below adapt these ideas to the level structures
and the characteristic hypotheses of Theorem~\ref{thm:main}.

Gaitsgory and Raskin establish the unramified restricted equivalence in
characteristic zero and develop geometric specialization and its
compatibility with the Hecke action
\cite[Main Theorem~1.3.9,\S4]{GR}. We use their characteristic-zero
equivalence to produce a compact eigencomplex, and prove the required
level-dependent specialization and nonvanishing assertions here.
The twisted nodal geometry belongs to the automorphic gluing framework
of Nadler and Yun \cite[\S\S2--3]{NYAutomorphicGluing}. Their gluing theorem
is a categorical construction; the nonvanishing of the particular
specialized eigencomplex used below is a separate argument.

In Betti sheaf theory, Nadler and Yun construct a spectral action with
level structure and prove local constancy of the moving Hecke operators
on the nilpotent category
\cite[Theorems~6.2.2,6.3.7]{NY}. At the marked point we use Gaitsgory's
nearby-cycle construction of central sheaves \cite{GaitsgoryCentral},
the local spectral geometry of Arkhipov and Bezrukavnikov \cite{AB},
and its universal monodromic form due to Dhillon and Taylor \cite{DT}.
The latter allows arbitrary enhancement monodromy before the central
trace calculation forces it to be unipotent.

\subsection{The two specializations and the proof strategy}
The proof first constructs an eigensheaf with Borel level in characteristic
zero, and then specializes it to the curve of Theorem~\ref{thm:main}.
The characteristic-zero construction itself uses a different specialization:
a degeneration to a separating node creates the level structure.

Start with a pointed complex curve and a dense parameter with unipotent
boundary monodromy on its punctured complement. Join two copies of the curve
at their marked points and smooth the resulting node. On the second
punctured component, choose a parameter
with inverse boundary monodromy. The compatible finite quotients of the two
parameters glue on twisted nodal models and lift to the smooth generic
curve. Their inverse limit is an unramified dense parameter there, to which
we apply the characteristic-zero unramified correspondence. Taking the
compact spectral skyscraper at this parameter gives a nonzero eigencomplex
that is locally constructible.

The special fiber introduces enhanced flags. Put $T=B/R_u(B)$ and let
$\cA^+$ be the stack of bundles with an $R_u(B)$-reduction at the marked
point. The map $q:\cA^+\to\cA$ is a $T$-torsor. For a suitable prescribed
stabilizer type at the twisted node, the bundle stack is
\[
 (\cA_1^+\times\cA_2^+)/T.
\]
Nearby cycles, pullback to the product, and restriction to one enhanced
factor give a nonzero eigencomplex on $\cA_1^+$. Taking a nonzero perverse
cohomology, applying $q_!$, and taking perverse cohomology once more produce
the ordinary-flag perverse eigensheaf. The arguments below establish
nonvanishing and preserve the eigenstructure at each of these steps.

To return to positive characteristic, lift the original pointed curve and
the compatible finite quotients of its parameter to mixed characteristic.
Apply the preceding construction on the geometric generic fiber. Its
geometric nearby cycles give the sheaf on $\cA$ in
Theorem~\ref{thm:main}. The two specializations are shown in
Figure~\ref{fig:construction}.

\begin{figure}[htbp]
\centering
\begin{tikzpicture}[
 box/.style={draw,rounded corners=2pt,align=center,text width=3.5cm,
             minimum height=1.2cm,inner sep=5pt,font=\small},
 arr/.style={-{Stealth[length=2mm]},thick},
 lab/.style={font=\footnotesize,align=center}]
 \node[box] (unram) {Unramified eigencomplex\\on a smooth curve};
 \node[box,right=1.1cm of unram] (enh) {Eigencomplex\\on $\cA_1^+$};
 \node[box,right=1.1cm of enh] (ord) {Perverse eigensheaf\\on $\cA_1$};
 \draw[arr] (unram) -- node[lab,above=8mm]{nodal nearby cycles\\and restriction} (enh);
 \draw[arr] (enh) -- node[lab,above=8mm]{perverse cohomology\\and $q_!$} (ord);
 \node[lab,above=18mm of enh] {Construction in characteristic zero};
 \node[box,below=1.2cm of ord] (positive) {Perverse eigensheaf\\on $\Bun_{G,B,x}(X)$};
 \draw[arr] (ord) -- node[lab,right]{mixed-characteristic\\nearby cycles} (positive);
\end{tikzpicture}
\caption{The first specialization creates the level structure; the second
changes the characteristic. The arrow from enhanced to ordinary flags takes
perverse cohomology both before and after $q_!$. Each specialization also
requires a separate nonvanishing argument. The arrows record operations on
sheaves, not morphisms between the displayed moduli stacks.}
\label{fig:construction}
\end{figure}

The first difficulty is nonvanishing. Nearby cycles on a nonproper bundle
stack can annihilate a nonzero object; at the node they must also be
nonzero on the prescribed stabilizer type. Uniformization places a nonzero
generic stalk in a bounded Hecke space proper over a reference bundle of
that type. Its extension shows that the support approaches the required
special fiber. Theorem~\ref{thm:specialization-detection} then proves that
nearby cycles cannot vanish.

The detection arguments have a common local calculation but different
conclusions. Over a field, a nonnilpotent covector admits a transverse
Hecke modification with a nonzero moving-leg component. The two-Hecke
calculation and projective recovery use the lisse eigenrelation to show
that a function with that differential has zero vanishing cycles, excluding
the covector from singular support. Over a trait, assume that nearby cycles
vanish. The same local calculation forces nearby cycles to vanish on
hypersurface cuts whose differentials avoid the nilpotent cone. The cone
dimension bound allows enough simultaneous cuts to retain a nonzero generic
stalk while making the support finite over the trait. Its nearby cycles
are a nonzero sum of stalks, giving the contradiction. The two-Hecke and
projective slicing methods follow
\cite[Sections~3--7]{OpenAIFullSupport2026}; we prove their versions for
ordinary and enhanced Borel level and the simultaneous torus quotient,
under \ref{hyp:form}--\ref{hyp:nilpotent}.

Preserving the Hecke system requires more than commuting nearby cycles
with a fixed-point operator. The comparison must retain all moving legs,
including collisions, and the maps expressing convolution and fusion.
In input frames, each successive Hecke correspondence is a product with
a fixed Grassmannian Schubert model. The nearby-cycle exterior-product
comparison applies even when earlier modifications have made the input
singular. Iterating it and using proper convolution gives the entire
coherent system in Section~\ref{sec:specialization}.

Perverse cohomology presents a different issue: general Hecke functors
are not assumed to be perverse exact. Choose generators of the punctured
surface group. A framing of the parameter gives a monodromy tuple
$y\in Y=\Gd^{2g}$, whose simultaneous-conjugation orbit is closed and free
because the parameter is dense and $\Gd$ is adjoint. For each representation
$V$, the equivariant bundle $Y\times V$ admits an equivariant frame after
inverting an
invariant function $f_V$ nonzero at $y$. Under the Betti spectral action,
this frame identifies the fixed-point Hecke functor with $\dim V$ copies
of the identity wherever $f_V$ acts invertibly. The eigencomplex belongs
to this subcategory because $f_V$ acts on it by the scalar $f_V(y)$.
Inverting these functions gives a category stable under perverse
truncation, on which all fixed-point Hecke functors are exact.
Nadler--Yun's local constancy theorem extends this exactness to moving legs
and their collisions. Functorial truncation then preserves the original
eigenmaps and all their compatibilities (Section~\ref{sec:perverse}).

Finally, $q_!$ can vanish on an enhanced-flag object with arbitrary torus
monodromy. A central-sheaf trace identity relates that monodromy to the
boundary monodromy of the parameter. Unipotence of the latter forces
unipotence along the enhancement fibers and hence nonvanishing of $q_!$.
To globalize the trace identity, we retain independent input and output
enhancements during collision at the marked point. After removing the
contractible real radial directions of the enhancement torus, the bounded
pushforward is proper. This comparison preserves the monodromy on each
central factor, as the trace requires (Section~\ref{sec:descent}).

The geometric nearby-cycle foundations are those of
Hansen--Scholze \cite{HansenScholze} and Gaitsgory--Raskin \cite{GR}.
Throughout, finite reductions may require different finite trait
extensions; we do not replace them by one extension supporting the entire
adic system. All sheaves and parameters are geometric, and no Frobenius
structure is required.

Section~\ref{sec:foundations} fixes the sheaf and Hecke conventions.
Sections~\ref{sec:geometry} and~\ref{sec:detection} establish the cotangent
geometry and detection arguments. After the Hecke comparison, perverse
cohomology, and torus descent of
Sections~\ref{sec:specialization}--\ref{sec:descent},
Section~\ref{sec:construction} constructs the two families and assembles
the proof of Theorem~\ref{thm:main}.
